\documentclass[12pt]{article}
\usepackage[utf8]{inputenc}
\usepackage[T1]{fontenc}
\usepackage{hyperref}
\title{Malyangapa Language Resources for K-12 Curriculum Development}
\author{Compiled by Claude Sonnet 5.5 (claude-sonnet-5-5) for Andr\'as Kornai}
\date{}
\begin{document}
\maketitle

\section{Introduction}
Malyangapa (ISO 639-3 \texttt{yga}) is an extinct Australian language of far north-west New South Wales, centred on Milparinka. Sources give 1976 as the year of the last known speaker's death. It belongs to the Yarli subgroup with Wadikali and Yardliyawara. The main record comes from Stephen Wurm's 1957 fieldwork with Hannah Quayle: a vocabulary of 357 items and 386 sentences. Peter Austin has worked on this material.

\section{Primary Educational Resources}
No school materials were found in the search results.

\section{Linguistic and Reference Materials}
\begin{itemize}
\item Peter K. Austin, Malyangapa project page: \url{https://peterkaustin.com/projects/malyangapa-project/}
\item Wikipedia, Yarli language: \url{https://en.wikipedia.org/wiki/Yarli_language}
\item ASJP word list for Malyangapa: \url{https://asjp.clld.org/languages/MALYANGAPA}
\item Academic paper (title not shown in the search result): \url{https://www.academia.edu/2099279}
\end{itemize}

\section{Religious and Community Resources}
\begin{itemize}
\item The Yarli Project (Corner Country, document 4): \url{https://visitcornercountry.com.au/wp-content/uploads/2024/03/The-Yarli-Project-4.pdf}
\end{itemize}
The search result showed only the title and URL of this document, so its contents are unverified.

\section{Digital and Audio Resources}
Austin's page notes that some Hercus and Beckett recordings had not yet been incorporated into the data files. No public audio resources were confirmed.

\section{Curriculum Development Context}
The language is extinct, so a K-12 programme would be heritage and revitalisation work, not language instruction. The 1957 Wurm materials are small, which limits what can be taught. Any programme would need to be led by Malyangapa descendants and local Aboriginal community bodies. This survey found no existing school programme.
\end{document}